\section[Locally exposed cycle rank]{A coarse system for the locally exposed cycles}
\label{sec:op3-local-cycle-rank}

\paragraph{Status and parameters.}
This section gives a \statusproved{} construction as a proof draft awaiting
independent review. It applies to arbitrary finite simple connected
unweighted graphs in the same local degree/adjacency and exact-real word
models. Its work depends on the cycle ranks of the final active support
and of the locally revealed graph. These dependencies remain explicit;
the result does not resolve graph-uniform OP3. The published height-bounded
top-tree algorithm remains an explicit \statussource{} import, and the
implemented audit rebuilds application hierarchies as reference work.

\begin{lemma}[Revealed cycle rank and exceptional incidences; \statusproved]
\label{lem:op3-revealed-cycle-rank}
For a nonempty connected active set $U$, let
\[
 r(U)=|E(G[U])|-|U|+1,\qquad
 a_j=|N(j)\cap U|\quad(j\in\partial U),
\]
and define the revealed cycle rank
\begin{equation}
 q(U)=r(U)+\sum_{j\in\partial U}(a_j-1).
 \label{eq:op3-revealed-cycle-rank}
\end{equation}
It is the cycle rank of the connected graph on $U\cup\partial U$ with
exactly the original edges incident to $U$. Both $r(U)$ and $q(U)$ are
nondecreasing under connected admissions, and
$r(U)\leq q(U)\leq\operatorname{vol}(U)$.
There are at most $q(U)-r(U)$ exceptional candidates with $a_j\geq2$,
and their total number of active parents is at most $2(q(U)-r(U))$.
All these quantities are maintained from scanned incidences only.
\end{lemma}
\begin{proof}
The revealed graph has $|U|+|\partial U|$ vertices and
$|E(G[U])|+\sum_j a_j$ edges, proving the identity. Admitting a candidate
with $a$ active parents increases $r$ by $a-1$ and removes the same amount
from the boundary sum. Each newly scanned incidence to a previously
discovered boundary vertex increases $q$ by one; an incidence introducing
a new boundary vertex changes its edge and vertex counts equally.
Thus both ranks are nondecreasing. Each exceptional summand $a_j-1$ is
at least one, and $a_j\leq2(a_j-1)$ for $a_j\geq2$. Finally the revealed
edge count is at most the number of scanned entries $\operatorname{vol}(U)$.
No edge between two still-inactive vertices is queried or included.
\end{proof}

\begin{lemma}[Two-port partition with permanent row homes; \statusproved]
\label{lem:op3-marked-tree-partition}
Let $T$ be an insertion spanning tree of $G[U]$, permanently rooted at $v$,
and let $R$ be the $r=r(U)$ remaining active edges. If $r\geq1$, there is
a retained set $P$ of at most $4r$ active vertices and a partition of
$E(T)$ into $|P|-1$ connected tree pieces, each with exactly two ports in
$P$. Distinct pieces have disjoint interiors and share only retained ports.
The total number of vertex copies over all pieces is
$|U|-1+(|P|-1)$.
The partition, original-label maps and permanent row-home component map
can be constructed in $O((|U|+r)\log(2+\operatorname{cvol}(U)))$ word
work from cached active tree edges, using comparison dictionaries.
An ordinary one-parent admission subsequently extends just one piece;
no partition change or global re-homing is required.
\end{lemma}
\begin{proof}
Mark $v$ and all endpoints of edges in $R$, giving a set $S$ of at most
$2r+1$ vertices. Let $K$ be the minimal subtree of $T$ connecting $S$.
Add every degree-at-least-three vertex of $K$ to $S$ to obtain $P$.
Every leaf of $K$ is marked, so the number of unmarked branching vertices
is at most $|S|-2$. Consequently $|P|\leq2|S|-2\leq4r$.
The maximal paths in $K$ between vertices of $P$ form a coarse tree with
$|P|-1$ edges. Since $v\in P$ and $T$ is rooted at $v$, the endpoints of
each such path are comparable; choose its rootward endpoint as its anchor.

Assign all unmarked branches attached to an interior path vertex to that
path's piece. At a retained vertex other than $v$, assign all its off-$K$
branches to its rootward path piece. At $v$, choose an incident piece for
any off-$K$ branches. The pieces are connected trees whose only shared
vertices are their two retained endpoints. Their edge sets partition $T$.
If there are $c$ pieces, summing the tree identity
$|V(B)|=|E(B)|+1$ gives total copy count $|U|-1+c$.

Home vertex $i\ne v$ on its permanent insertion-parent edge, and home $v$
on the first admitted edge. A row stays on that physical tree edge through
every repartition. A retained nonroot vertex therefore homes its row in
the rootward piece. At $v$, if its home edge is off $K$, assign that branch
to the selected home piece; otherwise its path already determines the
home piece. Each ordinary minimum row appears exactly once. Other copies
of a shared retained vertex contribute no duplicate reporter row.

To build $K$, process the cached insertion order backward, marking whether
each subtree contains a marked vertex. A tree edge is in $K$ exactly when
its farther subtree is marked. Count these incidences, identify $P$, and
traverse the resulting nonbranching paths once. Process the insertion order
forward to assign off-$K$ edges to their parent's home piece. This reads
$O(|U|+r)$ records and writes the same order of component, copy and home
records; dictionaries and deterministic sorting add the stated logarithm.
There is no ambient adjacency scan or reserved ambient vertex array.

If a new candidate has one active parent, it adds a leaf to $T$ and does
not change $R$, $S$, $K$ or $P$. Assign its new edge to its parent's home
piece. That piece contains the parent, even when the parent is shared.
The new vertex is unmarked, so both ports and every other piece remain
unchanged. Its own permanent home is its new insertion edge.
\end{proof}

\paragraph{The tree regime.}
When $r=0$, use one piece, the whole active tree, with its sole port $v$.
The ordinary reporter is the scalar tree reporter. There may still be many
exceptional inactive candidates, so the single-exception argument of the
unicyclic theorem must not be used here.

\begin{lemma}[Shared-port assembly and cycle restoration; \statusproved]
\label{lem:op3-coarse-cycle-assembly}
For every piece $B$, maintain its shifted Schur state
$(\boldsymbol H_B,\boldsymbol h'_B)$ at its one or two ports, using
original full degrees and
$\beta_i=\mathbf1_{i=v}-\lambda d_i-Cd_i$, where $C=1/\bar\alpha$.
Embed and sum these states on the retained coordinates $P$. If $i\in P$
appears in $m_i$ pieces, subtract $(m_i-1)d_i$ from its summed diagonal
and $(m_i-1)\beta_i$ from its summed load. For every extra edge
$\{i,j\}\in R$, subtract $\gamma$ from the two off-diagonal entries and
add $\gamma C$ to both endpoint loads. The resulting system is the exact
Schur system of the full active matrix and shifted load on $P$.
It is positive definite. Its solution supplies the exact port values for
every component reporter and named-coordinate query.
\end{lemma}
\begin{proof}
Before interior elimination, each piece contributes one $d_i$ and one
$\beta_i$ at every represented vertex, plus the unique contributions from
its own tree edges. The interiors are disjoint and have no edges to other
interiors, while all extra-edge endpoints are retained. Eliminating the
interiors can therefore be done independently within pieces. At shared
ports, the stated diagonal and load subtractions remove exactly the
duplicated fixed vertex terms. The extra-edge update is the per-edge
shifted rule from \cref{lem:op3-uniform-negative-shift}. Thus the assembled
system is precisely the full principal matrix's Schur complement, with its
correct shifted right-hand side. Positive definiteness follows from that
of $\boldsymbol M_{UU}$. In particular, an uncorrected piece or spanning-tree
response need not be positive in the physical coordinates; no obstacle
truncation is applied to these affine response objects.
\end{proof}

\paragraph{Local update and query algorithm.}
Cache each boundary candidate's active-parent list. A one-parent candidate
belongs to its parent's ordinary heap with key $\lambda d_j/\gamma-C$.
On acquiring a second parent, remove its ordinary row, refresh the old
parent's home edge, and put it in an exceptional dictionary. Later parents
only extend its incidence record; the candidate stays outside all ordinary
heaps. At every checkpoint, solve the coarse system and query every piece's
ordinary reporter at its resulting port values. Separately evaluate each
exceptional original gate
\[
 g_j=\gamma\sum_{i\in N(j)\cap U}u_i-\lambda d_j
\]
using \cref{lem:op3-cluster-point-query} in each named parent's home piece.
All quiet queries are included. Either type of strict positive gate can
supply the next admission; stop only when all queried signs are nonpositive.

An admission with $a\geq2$ active parents keeps one new edge in $T$ and
adds the other $a-1$ edges to $R$. Rebuild the marked partition and the
current component hierarchies from cached active edges. This occurs at
most $r(U_{\rm final})$ times. A one-parent admission instead uses one
source leaf insertion in the parent's home piece and the needed old home
payload updates. The source restores each affected piece's fixed one or
two external ports before any new gate query. Separate local structural
copies of shared ports are glued only by the coarse equation; they are not
additional vertices of the PageRank problem.

\begin{theorem}[A local bound in active and revealed cycle rank; \statusproved]
\label{thm:op3-local-cycle-rank}
The preceding deterministic algorithm returns the exact obstacle solution
on an arbitrary finite simple connected unweighted graph with at least two
vertices, for $0<\alpha<1$ and $\lambda>0$.
For a nonempty returned positive support $U$, put
$V=\operatorname{cvol}(U)$, $r=r(U)$, $q=q(U)$ and $L=\log(2+V)$.
Its total exact-real word work is
\begin{equation}
 O\!\left(1+V\bigl[(1+r)L^4+(1+r)^3+qL^2\bigr]\right),
 \label{eq:op3-cycle-rank-work}
\end{equation}
and its space, including all historical application versions, is
\begin{equation}
 O\!\left(1+(1+r)VL^3+(1+r)^2\right).
 \label{eq:op3-cycle-rank-space}
\end{equation}
An empty support is certified in constant work by $\lambda d_v\geq1$.
With $\lambda=\epsappr/2$ and $0<\epsappr<1$, it gives original ACL residual between zero and
$\lambda\boldsymbol d$ and fully charged local work
$\widetilde O((1+r^3+q)/\epsappr)$.
With $\lambda=\rho$, it gives the exact RPPR solution, of objective gap zero,
in $\widetilde O((1+r^3+q)/\rho)$ work in the nonzero regime.
\end{theorem}
\begin{proof}
The singleton and zero cases are as in \cref{thm:op3-local-unicyclic}.
By \cref{lem:op3-coarse-cycle-assembly}, every checkpoint has exact full
active-face coordinates at all retained ports. The component reporters
cover each ordinary heap minimum exactly once, while the exceptional
queries use every actual active parent. Thus every admission has a strict
positive original gate and the block update of \cref{thm:block-schur}
keeps every admitted coordinate positive and increases the old ones.
Termination gives the exact obstacle KKT conditions. All admitted rows
belong to the unique final positive support, giving at most $|U|$ iterations
and a single final recovery. The residual and mass arguments from the
unicyclic theorem apply unchanged.

There are $O(1+V)$ original row entries, degree replies, discovered-label
records and incidence insertions. Ordinary heap insertions and lazy
deletions have the same total count. A new second-parent incidence causes
at most one extra old home-payload refresh, so there are $O(1+V)$ such
refreshes in total, even if a single admitted row creates many exceptional
candidates. Source leaf insertions, exposure restoration and height-bounded
payload walks cost $O((1+V)L^4)$ work and allocate $O((1+V)L^3)$ persistent
application state. Component copies use local identifiers; dictionaries,
capacity changes and their data copies are charged as in
\cref{lem:op3-top-tree-payload}.

By \cref{lem:op3-marked-tree-partition}, a cycle-birth rebuild visits
$O(|U|+r)$ cached words and creates pieces of total size $O(|U|+r)$.
Initialize their source hierarchies in parent-before-child order. Since
$r\leq V$, one such rebuild costs $O(VL^4)$ work and $O(VL^3)$ newly
allocated application state. At most $r$ rebuilds give the first terms in
\eqref{eq:op3-cycle-rank-work} and \eqref{eq:op3-cycle-rank-space}.
Per-component capacity doubling between rebuilds has the same geometric-sum
bound. Permanent original home edges and depths do not change, but their
component map and structural leaf locators are rebuilt and charged.

Each coarse system has dimension $O(1+r)$. Zeroing its dense matrix, copying
it for elimination, assembling piece entries, correcting duplicates and
restoring extra edges cost $O((1+r)^2+(1+r)L)$ word work. A deterministic
dense positive-definite solve costs $O((1+r)^3)$ work and $O((1+r)^2)$
temporary words. There is one solve per checkpoint, including termination.
Querying all $O(1+r)$ component reporters costs $O((1+r)L^2)$ per checkpoint.
By \cref{lem:op3-revealed-cycle-rank}, all exceptional parent incidences
number at most $2q$, so their named queries and dictionary work cost
$O((1+q)L^2)$ per checkpoint, including every failed gate check and the
candidate list search. These bounds are absorbed in
\eqref{eq:op3-cycle-rank-work}, since $|U|\leq V$.
If incidence lists are copied as tuples, as in the exact driver, their total
copy work is at most $O((1+q)V)$, also covered by that bound; old temporary
tuple copies need not be retained.

Current component copies and their original-label maps use $O(V+r)$ words.
Final recovery traverses each current component once, checks equal values
at shared ports, adds $C$ and emits each original coordinate once. Its work
is $O((V+r)L)$. Historical root pointers and their solved one/two-port
values cost $O((1+r)|U|)$ words and are included in the allocated-state
bound. Temporary coarse matrices and query arithmetic are paid work but
need not be saved as historical application versions. Finally
$\lambda\operatorname{vol}(U)<1$ implies $V=O(1/\lambda)$, proving both
accuracy-namespace consequences.
\end{proof}

\paragraph{Relation to OP3.}
The support ranks are output parameters, not a supplied decomposition or an
ambient preprocessing budget. For fixed revealed cycle rank this reaches
OP3's $\widetilde O(1/\epsappr)$ ACL scale, while the general bound retains
the explicit rank factors. Separately, the sufficient structural condition
$1+r^3+q=O(\alpha^{-1/2})$ puts the exact regularized solve at OP2's
conjectured $\widetilde O(1/(\rho\sqrt\alpha))$ scale. These are distinct
accuracy namespaces; no semantic PPR conversion is inferred here.
Large active cycle rank leaves a cubic coarse-system factor, and many quiet
exceptional candidates leave repeated named-query work. Neither dependency
is proved necessary. Removing or separately paying them is \statusopen{}.

\paragraph{Measured implementation boundary.}
\texttt{local\_cycle\_rank\_clusters.py} implements local incidence discovery,
permanent home edges, cycle-birth partition changes, shared-port correction,
coarse solves, complete exceptional gates and one final output. Independent
dense Schur references check the assembled matrix and shifted load on every
audited positive face, in addition to full obstacle and original-residual
comparisons. The audit exercises multiple simultaneous exceptional candidates
and single admissions that create several cycles. Rebuilt component
hierarchies, their full parent indices and dense reference recovery are
separately counted audit work. The asymptotic theorem uses the same verified
published online hierarchy as the earlier tree and unicyclic constructions.
The sweep covers all 995 connected atlas graphs through seven vertices,
every seed, four parameter pairs and two policies: 54,240 executions.
Twenty larger explicit cases and five finite implicit private-star cases
extend the families. In total, 217,410 positive faces, 168,368 independently
assembled coarse Schur systems and 483,413 original gates are checked.
The atlas alone has 33,074 admissions creating multiple cycles and 166,405
exceptional gate checks, of which 3,441 are quiet. Retained component roots
are rechecked 198,001 times. The five implicit examples inspect only
58, 51, 48, 62 and 70 original entries, respectively, and never scan an
inactive hub row. Full parameters, counts and source hashes are in
\texttt{LOCAL\_CYCLE\_RANK\_CLUSTER\_AUDIT.json}.
